% -----------------------------------------------------------------------------
% Final inferential analysis across MKP, CEC2022, and HRES2--H2.
% This is an input fragment, not a standalone LaTeX document.
% Required packages: amsmath, booktabs, graphicx, and array.
%
% Data sources:
% - MKP: resultados finales/run_{dtw,ddtw}{,_3k}_mkp/
% - CEC2022: resultados finales/run_dtw2022_job_19070/ and
%            resultados finales/run_cec2022_ddtw_job_19066/
% - HRES2--H2: resultados finales/run_{dtw,ddtw}_hres2/
% -----------------------------------------------------------------------------

\section{Cross-Domain Inferential Statistical Analysis}
\label{sec:statistical-analysis-final}

This section evaluates whether the proposed DTW- and DDTW-driven hybrid
frameworks differ significantly from the standalone metaheuristics in the
three experimental domains: the Multidimensional Knapsack Problem (MKP), the
IEEE CEC2022 continuous benchmark, and the HRES2--H2 engineering case. Each
method was evaluated over 31 stochastic runs. In MKP and CEC2022, corresponding
run indices followed the same seed schedule and were treated as matched
experimental blocks. In HRES2--H2, this pairing is documented only for the
direct DTW--DDTW comparison; the standalone-method samples were therefore
treated as independent.

\subsection{Statistical protocol}

Normality was assessed with the Shapiro--Wilk test at $\alpha=0.05$. Since
normality was rejected in a substantial fraction of the framework
distributions, inferential conclusions were based on non-parametric tests. For
the matched MKP and CEC2022 observations, the Friedman test was used as the
omnibus comparison and the two-sided Wilcoxon signed-rank test was used for
pairwise contrasts. The HRES2--H2 hybrid--baseline comparisons used the
two-sided Mann--Whitney $U$ test because the archived baseline samples are
independent rather than matched. Family-wise error was controlled with Holm's
step-down procedure: over the six MKP configuration contrasts, over the seven
hybrid--baseline tests within each MKP instance, over the five
hybrid--baseline tests within each CEC2022 function, over the 12 direct
CEC2022 DTW--DDTW contrasts, and separately over the ten HRES2--H2 comparisons
of each framework variant. A comparison was considered significant when
$p_{\mathrm{Holm}}<0.05$.

The direction of every pairwise result is reported from the perspective of
the hybrid framework. Thus, W/T/L denotes a significant hybrid win, a
non-significant tie, or a significant hybrid loss, respectively. Larger
fitness is better in MKP, whereas smaller objective values are better in
CEC2022 and HRES2--H2. Raw objective values from different MKP instances or
CEC2022 functions were not pooled because they have different scales. Global
conclusions instead use normalized gaps, within-problem ranks, and aggregated
within-problem pairwise decisions.

\begin{table}[htbp]
\centering
\caption{Statistical design and principal comparison against standalone
metaheuristics in each domain.}
\label{tab:final-stat-protocol}
\small
\resizebox{\textwidth}{!}{%
\begin{tabular}{lllll}
\toprule
Domain & Performance metric & Standalone methods & Multiplicity family & Main outcome \\
\midrule
MKP & Normalized BKS gap & 7 methods & Wilcoxon--Holm: 7 per instance & Portfolio advantage, not universal dominance \\
CEC2022 & Objective value/rank & 5 methods & Wilcoxon--Holm: 5 per function & Best aggregate rank in both campaigns \\
HRES2--H2 & LCOE (CNY/kWh) & 10 methods & Mann--Whitney--Holm: 10 per variant & Higher baseline means and significant distributional differences \\
\bottomrule
\end{tabular}%
}
\end{table}

\subsection{Multidimensional Knapsack Problem}
\label{sec:final-stat-mkp}

The MKP analysis covered the first instance of each of the nine Chu--Beasley
families and four framework configurations: DTW--1K, DDTW--1K, DTW--3K, and
DDTW--3K. Shapiro--Wilk rejected normality in 14 of the 36
instance--configuration distributions. Therefore, the comparison among the
four configurations used the mean percentage gap to the best known solution
(BKS), with lower values indicating better performance.

The Friedman test detected a significant difference among the four
configurations ($\chi_F^2=24.0667$, $p=2.419\times10^{-5}$). DTW--3K obtained
the best mean rank (1.22), followed by DDTW--3K (1.78), DTW--1K (3.11), and
DDTW--1K (3.89). All six paired configuration comparisons remained
significant after Holm correction (Table~\ref{tab:final-mkp-config-tests}).
Increasing the budget from 1K to 3K reduced the average gap from 2.237\% to
1.805\% for DTW (19.3\% relative reduction) and from 2.391\% to 1.876\% for
DDTW (21.5\% relative reduction).

\begin{table}[htbp]
\centering
\caption{Paired Wilcoxon comparisons of the mean MKP gap across the nine
instances. Holm adjustment was applied to the six contrasts.}
\label{tab:final-mkp-config-tests}
\begin{tabular}{lrr}
\toprule
Configuration contrast & Raw $p$ & $p_{\mathrm{Holm}}$ \\
\midrule
DTW--1K vs. DDTW--1K   & 0.027344 & 0.039063 \\
DTW--1K vs. DTW--3K    & 0.003906 & 0.023438 \\
DTW--1K vs. DDTW--3K   & 0.003906 & 0.023438 \\
DDTW--1K vs. DTW--3K   & 0.003906 & 0.023438 \\
DDTW--1K vs. DDTW--3K  & 0.003906 & 0.023438 \\
DTW--3K vs. DDTW--3K   & 0.019531 & 0.039063 \\
\bottomrule
\end{tabular}
\end{table}

Within each instance and campaign, the hybrid framework was compared with
PSO, GWO, WOA, EHO, ACO, GA, and SA. All 36 per-instance Friedman tests were
significant. Table~\ref{tab:final-mkp-wtl} aggregates the Holm-adjusted
decisions without mixing raw fitness values across instances.

\begin{table}[htbp]
\centering
\caption{MKP wins/ties/losses of each hybrid configuration against the seven
standalone metaheuristics over nine instances. Each cell is W/T/L from the
hybrid perspective.}
\label{tab:final-mkp-wtl}
\small
\resizebox{\textwidth}{!}{%
\begin{tabular}{lrrrrrrrr}
\toprule
Configuration & PSO & GWO & WOA & EHO & ACO & GA & SA & Total \\
\midrule
DTW--1K  & 9/0/0 & 5/4/0 & 6/3/0 & 8/1/0 & 7/2/0 & 5/1/3 & 1/8/0 & \textbf{41/19/3} \\
DDTW--1K & 9/0/0 & 5/4/0 & 5/4/0 & 9/0/0 & 7/2/0 & 4/2/3 & 1/4/4 & \textbf{40/16/7} \\
DTW--3K  & 9/0/0 & 5/3/1 & 5/3/1 & 9/0/0 & 7/2/0 & 5/1/3 & 1/6/2 & \textbf{41/15/7} \\
DDTW--3K & 9/0/0 & 5/3/1 & 5/3/1 & 9/0/0 & 7/2/0 & 4/1/4 & 1/6/2 & \textbf{40/15/8} \\
\bottomrule
\end{tabular}%
}
\end{table}

The hybrids were particularly consistent against PSO, EHO, and ACO. GA was
the strongest competitor, producing significant advantages in three
instances against DTW--1K, three against DDTW--1K, three against DTW--3K, and
four against DDTW--3K. SA was frequently statistically tied with the hybrid
and was significantly better in selected 1K and 3K campaigns. Consequently,
the MKP evidence supports robust portfolio performance, while it does not
support a claim of universal superiority over every standalone method.

\subsection{IEEE CEC2022 continuous benchmark}
\label{sec:final-stat-cec2022}

The final CEC2022 analysis uses the completed DTW job 19070 and DDTW job
19066. Each campaign compared its hybrid control with PSO, GWO, WOA, EHO, and
ACO on functions F1--F12 at $D=10$. Shapiro--Wilk rejected normality for the
hybrid in 10 of the 12 functions in each campaign. All 12 per-function
Friedman tests were significant for both DTW and DDTW, demonstrating that the
six algorithms did not have equivalent performance distributions on any
function.

Ranks were calculated within each function and then averaged across the 12
functions. The hybrid obtained the best aggregate rank in both campaigns:
2.323 for DTW and 2.358 for DDTW (Table~\ref{tab:final-cec-ranks}). This
within-function aggregation avoids distortion from the different CEC2022
objective scales and bias terms.

\begin{table}[htbp]
\centering
\caption{Aggregate mean ranks across the 12 CEC2022 functions. Lower ranks
indicate better performance.}
\label{tab:final-cec-ranks}
\begin{tabular}{lrr}
\toprule
Algorithm & DTW campaign & DDTW campaign \\
\midrule
Hybrid framework & \textbf{2.323} & \textbf{2.358} \\
EHO              & 3.156 & 3.142 \\
ACO              & 3.257 & 3.269 \\
PSO              & 3.372 & 3.364 \\
GWO              & 3.616 & 3.598 \\
WOA              & 5.277 & 5.269 \\
\bottomrule
\end{tabular}
\end{table}

The Holm-adjusted pairwise outcomes are summarized in
Table~\ref{tab:final-cec-wtl}. DTW achieved 38 significant wins, 16 ties, and
6 losses in 60 comparisons, whereas DDTW achieved 37 wins, 19 ties, and 4
losses. WOA showed the clearest framework advantage, with 11 wins and one tie
for each alignment variant. ACO was the most competitive baseline, accounting
for three DTW losses and two DDTW losses.

\begin{table}[htbp]
\centering
\caption{CEC2022 Holm-adjusted Wilcoxon outcomes against the standalone
metaheuristics. Counts are W/T/L from the hybrid perspective over 12
functions.}
\label{tab:final-cec-wtl}
\begin{tabular}{lrr}
\toprule
Comparator & DTW & DDTW \\
\midrule
PSO & 7/3/2  & 6/5/1 \\
GWO & 8/4/0  & 8/4/0 \\
WOA & 11/1/0 & 11/1/0 \\
EHO & 7/4/1  & 6/5/1 \\
ACO & 5/4/3  & 6/4/2 \\
\midrule
Total & \textbf{38/16/6} & \textbf{37/19/4} \\
\bottomrule
\end{tabular}
\end{table}

The direct paired DTW--DDTW comparison was also conducted independently for
each function and adjusted across the 12 tests. None of the contrasts remained
significant after Holm correction ($0/12$). DTW had the lower mean at full
precision on seven functions and DDTW on five, but this descriptive difference
does not establish statistical superiority of either alignment strategy.
Therefore, the CEC2022 results support broad competitiveness against the
standalone methods and a landscape-dependent balance between DTW and DDTW.

\subsection{HRES2--H2 engineering case}
\label{sec:final-stat-hres2}

The engineering experiment compared each hybrid with PSO, GWO, WOA, EHO,
ACO, ABC, ILS, SA, TS, and VNS using the levelized cost of electricity (LCOE)
as the common minimization metric. Each objective evaluation simulated 8,760
hourly dispatch steps over the same fixed synthetic weather profile. The
Shapiro--Wilk test rejected normality for both hybrid LCOE distributions
($p<10^{-11}$).

The HRES2--H2 archive does not preserve the individual baseline observations
or a common-seed correspondence between each baseline and the hybrid. Thus,
the archived row-wise Friedman and Wilcoxon calculations are not used as
inferential evidence here. Instead, Table~\ref{tab:final-hres-tests} reports
the archived mean LCOE and the archived two-sided Mann--Whitney $p$-values
after a reproducible Holm adjustment over the ten comparisons in each
campaign. Every baseline had a higher mean LCOE than its hybrid control, and
all 20 distributional contrasts remained significant after adjustment. This
yields a 10/0/0 win/tie/loss record for each framework variant.

\begin{table}[htbp]
\centering
\caption{HRES2--H2 comparison against standalone metaheuristics. The reported
$p$-values are two-sided Mann--Whitney values adjusted by Holm over the ten
independent-sample comparisons of each framework variant.}
\label{tab:final-hres-tests}
\scriptsize
\resizebox{\textwidth}{!}{%
\begin{tabular}{lrrrr}
\toprule
& \multicolumn{2}{c}{DTW campaign} & \multicolumn{2}{c}{DDTW campaign} \\
\cmidrule(lr){2-3}\cmidrule(lr){4-5}
Method & Mean LCOE & $p_{\mathrm{Holm}}^{\mathrm{MW}}$ & Mean LCOE & $p_{\mathrm{Holm}}^{\mathrm{MW}}$ \\
\midrule
Hybrid control & \textbf{0.267160} & -- & \textbf{0.267160} & -- \\
GWO & 0.268066 & $6.384\times10^{-10}$ & 0.268395 & $2.272\times10^{-9}$ \\
VNS & 0.267806 & $2.232\times10^{-10}$ & 0.267806 & $1.090\times10^{-9}$ \\
TS  & 0.267188 & $1.242\times10^{-10}$ & 0.267184 & $6.880\times10^{-10}$ \\
ILS & 0.267184 & $1.242\times10^{-10}$ & 0.267183 & $6.880\times10^{-10}$ \\
ACO & 0.273288 & $7.639\times10^{-7}$  & 0.273617 & $1.316\times10^{-5}$ \\
PSO & 0.277289 & $1.285\times10^{-2}$  & 0.280068 & $1.573\times10^{-3}$ \\
WOA & 0.277227 & $2.373\times10^{-5}$  & 0.276477 & $1.797\times10^{-2}$ \\
ABC & 0.272971 & $8.591\times10^{-11}$ & 0.270993 & $1.329\times10^{-10}$ \\
EHO & 0.277728 & $4.333\times10^{-7}$  & 0.278524 & $1.431\times10^{-8}$ \\
SA  & 0.286431 & $8.591\times10^{-11}$ & 0.290062 & $9.340\times10^{-11}$ \\
\bottomrule
\end{tabular}%
}
\end{table}

The archive stores per-run LCOE values to six decimal places. On these values,
the paired same-seed Wilcoxon test did not detect a difference between DTW and
DDTW ($p=0.593$). The mean difference was only $7.42\times10^{-7}$ CNY/kWh,
below the reporting precision.
Accordingly, the experiment provides no evidence of an LCOE difference between
the variants; this non-significant result is not, by itself, a formal
equivalence test. The conclusion is restricted to the fixed synthetic weather
profile and to LCOE. Moreover, the standalone-method archive retains summary
statistics and test outputs, but not the individual baseline LCOE observations
or per-run LCOH, hydrogen-production, feasibility, and component-sizing
records.

\subsection{Cross-domain interpretation}
\label{sec:final-stat-synthesis}

The three domains provide complementary evidence. In MKP, the 3K budget
significantly improves solution quality and DTW--3K has the best configuration
rank, although GA and SA remain competitive on selected instances. In
CEC2022, each hybrid has the best aggregate rank in its own campaign and wins
most pairwise comparisons, but neither DTW nor DDTW is significantly superior
to the other. In HRES2--H2, both hybrids have lower mean LCOE than all ten
standalone methods, with significant independent-sample distributional
differences after Holm correction, while the paired test detects no
DTW--DDTW difference.
Overall, the results support the robustness of the adaptive portfolio across
discrete, continuous, and engineering optimization, without implying universal
dominance on every individual problem or formal equivalence of DTW and DDTW.
